

 title\+: Open\+F\+O\+AM as F\+MU permalink\+: doc\+\_\+\+Ofas\+F\+M\+U.\+html keywords\+: design sidebar\+: doc\+\_\+sidebar \subsection*{folder\+: doc }

\subsection*{Generate F\+MU from an Open\+F\+O\+AM case}


\begin{DoxyCode}
pip install fmu4foam
\end{DoxyCode}


creating an F\+MU from Open\+F\+O\+AM is simple and similar to the libary \href{https://github.com/NTNU-IHB/PythonFMU}{\tt Python\+F\+MU} which is achieved with just one line\+:


\begin{DoxyCode}
fmu4foam build -f TempControl.py -of OfCase/ --no-external-tool
\end{DoxyCode}


\subsubsection*{Open\+F\+O\+AM modifications}

As in the previous method, we extend Open\+F\+O\+AM with the additional libraries and add the following lines to the system/control\+Dict\+:


\begin{DoxyCode}
\textcolor{comment}{// loads addtional input and output function/class}
\textcolor{comment}{// e.g. new boundary conditions.}
libs(externalComm); 

\textcolor{comment}{// includes the file system/simulationParameters in the controlDict}
\textcolor{comment}{// and enables the definition of Variables marked with a $ }
\textcolor{comment}{// only look in the current directory and is short for:}
\textcolor{comment}{// #include "<case>/system/simulationParameters"}
\textcolor{preprocessor}{#include "simulationParameters"} 

functions
\{
    FMUController
    \{
        type            FMUController; \textcolor{comment}{// execute functionObjects FMUController}
        libs            (FMUController); \textcolor{comment}{// load "libFMUController.so"}
        host            $host; \textcolor{comment}{// connect to ip specified in simulationParameters}
        port            $port; \textcolor{comment}{// connect to port specified in simulationParameters}

    \}
\}
\end{DoxyCode}


the control\+Dict specifies two variables\+: {\itshape host} and {\itshape port} that need to be specified in the system/simulation\+Parameters\+:


\begin{DoxyCode}
\textcolor{comment}{/*--------------------------------*- C++ -*----------------------------------*\(\backslash\)}
\textcolor{comment}{| =========                 |                                                 |}
\textcolor{comment}{| \(\backslash\)\(\backslash\)      /  F ield         | OpenFOAM: The Open Source CFD Toolbox           |}
\textcolor{comment}{|  \(\backslash\)\(\backslash\)    /   O peration     | Version:  v2012                                 |}
\textcolor{comment}{|   \(\backslash\)\(\backslash\)  /    A nd           | Website:  www.openfoam.com                      |}
\textcolor{comment}{|    \(\backslash\)\(\backslash\)/     M anipulation  |                                                 |}
\textcolor{comment}{\(\backslash\)*---------------------------------------------------------------------------*/}
FoamFile
\{
    version     2.0;
    format      ascii;
    \textcolor{keyword}{class       }dictionary;
    location    \textcolor{stringliteral}{"system"};
    \textcolor{keywordtype}{object}      simulationParameters;
\}
\textcolor{comment}{// * * * * * * * * * * * * * * * * * * * * * * * * * * * * * * * * * * * * * //}

host            \textcolor{stringliteral}{"127.0.0.1"}; \textcolor{comment}{// defines host }
port            8000; \textcolor{comment}{// defines port}
STDGRAD         Gauss linear;
STDLAP          Gauss linear uncorrected;
writeInterval   0.1;


\textcolor{comment}{// ************************************************************************* //}
\end{DoxyCode}
 The host and port variable are specified in the F\+MU generation base class.


\begin{DoxyItemize}
\item host -\/$>$ connect to other computer
\item port -\/$>$ port on that computer
\item output\+Path -\/$>$ extract dir of the Open\+F\+O\+AM Case
\item oscmd -\/$>$ name of the system command (enables exectuion on docker or wsl for windows users in the future)
\item arguments -\/$>$ additional parameters for the Allrun script e.\+g. source environment in script
\end{DoxyItemize}

N\+O\+TE\+: The host variable needs to be quoted. To achieve this, we define it with two qoutes \textquotesingle{}\char`\"{}\+Name of the Variable\char`\"{}\textquotesingle{}. The other variables require no specific highlighting.


\begin{DoxyCode}
class OF2Fmu(Fmi2Slave):


    def \_\_init\_\_(self, **kwargs):
        super().\_\_init\_\_(**kwargs)
        #local host aka this computer
        self.host = '"127.0.0.1"' # Note the double quotes
        self.port = 8000 # port see above
        self.outputPath = "OFCase" # extract path relative to FMU
        self.oscmd = "bash -i" # system command to execute the Allrun 
        # additional parameters for the Allrun script
        # of2012 to source the enviroment in the Allrun script
        self.arguments = "" 

        self.register\_variable(String("host",
            causality=Fmi2Causality.parameter, variability=Fmi2Variability.tunable))
        self.register\_variable(Integer("port",
             causality=Fmi2Causality.parameter, variability=Fmi2Variability.tunable))
        self.register\_variable(String("outputPath", 
            causality=Fmi2Causality.parameter, variability=Fmi2Variability.tunable))
        self.register\_variable(String("oscmd", 
            causality=Fmi2Causality.parameter, variability=Fmi2Variability.tunable))
        self.register\_variable(String("arguments", 
            causality=Fmi2Causality.parameter, variability=Fmi2Variability.tunable))
\end{DoxyCode}


\subsubsection*{F\+MU generation input file}

The last part in the generation of the F\+MU is a python file specified after -\/f Temp\+Control.\+py\+: 
\begin{DoxyCode}
from pythonfmu.variables import Boolean
from FMU4FOAM import Fmi2Causality, Fmi2Variability, Real, Boolean , Integer, String
from FMU4FOAM import OF2Fmu


class TempControl(OF2Fmu):

    author = "John Doe"
    description = "A simple description"

    def \_\_init\_\_(self, **kwargs):
        super().\_\_init\_\_(**kwargs)

        # Input Output variable
        self.Qin = 0.0
        self.dTout = 0.0
        self.Tout = 298.0
        self.host = '"192.168.1.200"' # overwrite the ip
        self.outputPath = "Testing" # will extract the OfCase files in Testing

        self.register\_variable(Real("Qin", causality=Fmi2Causality.input))
        self.register\_variable(Real("dTout", causality=Fmi2Causality.output))
        self.register\_variable(Real("Tout", causality=Fmi2Causality.output))

        # Parameters
        self.writeInterval = 1.0 # changes outputInterval of the solver
        self.STDLAP = "Gauss linear corrected" # possiblity to change discreization
        self.STDGRAD = "pointCellsLeastSquares" # possiblity to change discreization


        self.register\_variable(String("STDGRAD", 
            causality=Fmi2Causality.parameter, variability=Fmi2Variability.tunable))
        self.register\_variable(String("STDLAP", 
            causality=Fmi2Causality.parameter, variability=Fmi2Variability.tunable))
        self.register\_variable(Real("writeInterval", 
            causality=Fmi2Causality.parameter, variability=Fmi2Variability.tunable))
\end{DoxyCode}


The F\+MU standard defines three relevant causalities\+:


\begin{DoxyItemize}
\item input -\/$>$ input parameters that change in time (assume to be R\+E\+AL aka floats)
\item output -\/$>$ output parameters that change in time (assume to be R\+E\+AL aka floats)
\item parameter -\/$>$ set at the start of the simulation and does not change (strings, float, interges, boolean)
\end{DoxyItemize}

N\+O\+TE\+: F\+MI standard 2.\+0 only allows for scalar data transfer but does not support data fields so an Open\+F\+O\+AM vector would be 3 input variables (This will change in F\+MI 3.\+0)

In the example above we specify following causalities\+:


\begin{DoxyItemize}
\item input\+:
\begin{DoxyItemize}
\item Qin
\end{DoxyItemize}
\item input\+:
\begin{DoxyItemize}
\item Tout
\item d\+Tout
\end{DoxyItemize}
\item parameter
\begin{DoxyItemize}
\item write\+Interval
\item S\+T\+D\+L\+AP
\item S\+T\+D\+G\+R\+AD
\item host (from base class)
\item port (from base class)
\item output\+Path (from base class)
\item oscmd (from base class)
\item arguments (from base class)
\end{DoxyItemize}
\end{DoxyItemize}

These parameters can be modified before running the F\+MU for example with the \href{https://github.com/CATIA-Systems/FMPy}{\tt fmpy} by calling


\begin{DoxyCode}
python -m fmpy.gui TempControl.fmu
\end{DoxyCode}


\{\% include image.\+html file=\char`\"{}fmpy-\/\+Temp\+Control.\+png\char`\"{} alt=\char`\"{}fmpy-\/gui\char`\"{} caption=\char`\"{}\+Display in Fmpy G\+U\+I\char`\"{} \%\} 